\documentclass[10pt]{article}
\usepackage[utf8]{inputenc}
\usepackage{url}
\usepackage{hyperref}
\usepackage{amsmath}
\usepackage{amsfonts}
\usepackage{amssymb}
\usepackage{graphicx}
\usepackage{float}
\usepackage{lipsum}
\usepackage{multicol}
\usepackage{xcolor}
\usepackage[super]{natbib}
\usepackage{physics}
\usepackage[font=small]{caption}
\DeclareMathOperator*{\argmin}{arg\,min}
\addtolength{\abovecaptionskip}{-3mm}
\addtolength{\textfloatsep}{-5mm}
\setlength\columnsep{20pt}

\usepackage[a4paper,left=1.50cm, right=1.50cm, top=1.50cm, bottom=1.50cm]{geometry}


\author{Hrishikesh Belagali}

\title{Adiabatic Quantum Computation of QUBO Problems}

\begin{document}
	
	\begin{center}
		{\Large \textbf{Adiabatic Quantum Computation of QUBO Problems}}\\
		\vspace{1em}
		{\large Hrishikesh Belagali}\\
		\vspace{1em}
		\textit{CMSE 890 Project}
	\end{center}

	\begin{center}
		\rule{150mm}{0.2mm}
	\end{center}		

	\vspace{5mm}
	
\begin{multicols*}{2}

\section{Introduction}
Combinatorial optimization problems are ubiquitous across fields, 
including logistics, finance, machine learning, physics, and many others. These problems
involve finding an optimal solution to a problem from a finite set of possible solutions,
which often grows rapidly with the size of the problem. Many combinatorial optimization problems 
are NP-hard, meaning that there is no known polynomial-time algorithm to solve them exactly.\\\\
Existing methods for solving combinatorial optimization problems include exact algorithms such as 
branch-and-bound\cite{AutomaticMethodSolving}, cutting planes\cite{CuttingPlanes1980}, dynamic programming\cite{bellman1957dynamic}, and systematic 
enumeration\cite{balasAdditiveAlgorithmSolving1965}. However, these methods can be computationally expensive and may not be feasible for large instances. 
Another class of methods are heuristic algorithms, such as simulated annealing\cite{kirkpatrickOptimizationSimulatedAnnealing1983}, tabu search\cite{gloverTabuSearch1993}, genetic algorithms\cite{sastryGeneticAlgorithms2005}, and many others. These algorithms are often used when exact methods are too slow, but they do not guarantee optimal solutions. \\\\
Quantum computing has been explored as a potential approach to solving combinatorial optimization problems more efficiently than classical algorithms, leveraging phenomena such as superposition and entanglement. 
One of the earliest algorithms proposed was quantum stochastic optimization, which was a procedure based on the quantum tunnel effect to escape local minima in optimization problems\cite{apolloniQuantumStochasticOptimization1989a}. It was later renamed to quantum annealing (QA) and proposed 
as an alternative to simulated annealing for solving combinatorial optimization problems, in which 
quantum fluctuations were used to escape local minima instead of thermal fluctuations\cite{PDFNumericalImplementation}. It was later shown that quantum annealing demonstrated faster convergence 
compared to simulated annealing for certain classes of problems\cite{kadowakiQuantumAnnealingTransverse1998}.
The experimental implementation of QA in a disordered quantum ferromagnet allowed for the algorithm to be 
considered as an alternative to the standard gate model, involving a designated quantum device\cite{QuantumAnnealingDisordered}. \\\\
Adiabatic quantum computation (AQC) is another optimization method which aims to leverage quantum mechanics to solve optimization problems. It relies on the adiabatic theorem of quantum mechanics, and works by encoding the solution to the problem in the ground state of a Hamiltonian. Then the system is slowly evolved from an initial Hamiltonian with a known ground state to the problem Hamiltonian. If the evolution is slow enough, the system will remain in the ground state, yielding the optimal solution to the problem. \\\\ 
It has been shown that non-stoquastic AQC, i.e. AQC with a non-stoquastic Hamiltonian is as powerful as the 
circuit model of quantum computation, and thus can solve any problem that a universal quantum computer can with atmost polynomial overhead\cite{AdiabaticQuantumComputation}.  \\\\
In this project, we will take a specific class of combinatorial optimization problems,
known as Quadratic Unconstrained Binary Optimization (QUBO) problems, and implement 
an adiabatic evolution algorithm to solve them. We will compare the results with classical heuristic 
methods and analyze the performance. The goal of the project is to understand the factors
that influence AQC's performance and behavior.
\section{Methods}
\subsection{Quadratic Unconstrained Binary Optimization}
The Quadratic Unconstrained Binary Optimization (QUBO) problem is a combinatorial optimization 
problem that minimizes a quadratic objective function defined over binary variables. 
Formally, given a symmetric matrix $\mathbf{Q} \in \mathbb{R}^{n \times n}$,
the QUBO problem can be expressed as:
$$\mathbf{x}^* = \argmin_{\mathbf{x} \in \{0,1\}^n} \mathbf{x}^\top \mathbf{Q} \mathbf{x}$$
Note that $\mathbf{Q}$ is generally assumed to be symmetric, however this is not a strict requirement since we can always symmetrize it by replacing $\mathbf{Q}$ with $\frac{\mathbf{Q} + \mathbf{Q}^\top}{2}$ without changing the optimization problem.
Several combinatorial problems can be reformulated as QUBO problems, including the Max-Cut problem,
Graph Coloring, the Traveling Salesman Problem, and many others~\cite{gloverTutorialFormulatingUsing2019}.
In this project, we will formulate the QUBO problem as a Hamiltonian and use adiabatic evolution to 
find the optimal solution. 
\subsection{Hamiltonian Formulation}
The QUBO problem can be written as 
$$f(\mathbf{x}) = \sum_{i=1}^n \sum_{j=1}^n \mathbf{Q}_{ij} x_i x_j$$ 
First, we will map the variables from $\{0,1\}$ to $\{-1,1\}$ using the 
transformation $$x_i = \frac{1 - s_i}{2}$$ where $s_i \in \{-1,1\}$ is a spin variable.
Then, $f$ can be rewritten as 
\begin{align*}
f(\mathbf{s}) &= \sum_{i=1}^n \sum_{j=1}^n \mathbf{Q}_{ij} \frac{1 - s_i}{2} \frac{1 - s_j}{2} \\ 
&= \sum_{i=1}^n\sum_{j=1}^n \mathbf{J}_{ij} s_i s_j + \sum_{i=1}^n \mathbf{h}_i s_i + c
\end{align*}
where $J_{ij}, h_i, c$ are derived from $\mathbf{Q}$. The full derivation is given in Appendix~\ref{sec:qubo_to_hamiltonian}. Note that the constant does not affect 
the optimization, so we can ignore it. The Hamiltonian is then defined by replacing the 
spin variables with the corresponding Pauli-Z operators.
$$\mathcal{H} = \sum_{i=1}^n\sum_{j=1}^n \mathbf{J}_{ij} Z_i Z_j + \sum_{i=1}^n \mathbf{h}_i Z_i$$
Now, the ground state of $\mathcal{H}$ corresponds to the optimal solution of the QUBO problem. It is 
found by implementing the adiabatic evolution as a time evolution operator, and measuring the final state in 
the computational basis.
\subsection{Adiabatic Evolution}
In adiabatic quantum computation, we start with an initial Hamiltonian $\mathcal{H}_0$ whose 
ground state is trivial to prepare. Then, we define a time-dependent Hamiltonian that 
interpolates between $\mathcal{H}_0$ and the problem Hamiltonian $\mathcal{H}$. The 
solution to the problem is obtained by evolving the system according to the time-dependent 
Schrodinger equation and measuring the final state. The key constraint is that the evolution
must be slow enough to ensure that the system remains in its ground state throughout the 
process, which is guaranteed by the adiabatic theorem. The total evolution time $T$ must satisfy 
$$T \gg \frac{1}{\Delta^2}$$
where $\Delta$ is the minimum energy gap between the ground state and the first excited state of the Hamiltonian during the evolution.
Here, we define the initial Hamiltonian as $$\mathcal{H}_0 = -\sum_{i=1}^n X_i$$
where $X_i$ is the Pauli-X operator acting on the $i$-th qubit. 
The ground state of $\mathcal{H}_0$ is $\ket{+}^{\otimes n}$, which is a uniform superposition of all possible states.
Next, we define the problem Hamiltonian $\mathcal{H}_Q$ as derived in the previous section.
Finally, we define the time-dependent Hamiltonian $\mathcal{H}_A(t)$ as 
$$\mathcal{H}_A(t) = (1 - s(t)) \mathcal{H}_0 + s(t) \mathcal{H}_Q$$
where $s(t)$ is a monotonic function that satisfies $s(0) = 0$ and $s(T) = 1$.
The evolution is Trotterized\ref{sec:trotterization} into discrete time steps, and at each step, 
we apply the corresponding unitary operator to evolve the state. Note that as $\Delta t \to 0$, 
the Hamiltonian is effectively constant during each time step, and the evolution can be 
approximated by a sequence of unitary operations corresponding to the Hamiltonian at each time step.
\subsection{Trotterization}
The time evolution operator for a time-dependent Hamiltonian $\mathcal{H}(t)$ is given by the time-ordered exponential:
$$U(T) = \mathcal{T} \exp\left(-i \int_0^T \mathcal{H}(t) dt\right)$$
However, this operator is difficult to compute directly. Instead, we can approximate it using Trotterization,
which breaks the evolution into small time steps $\Delta t$ and approximates the evolution as a product of exponentials of the Hamiltonian at each time step. \\\\
Let $T$ be the total evolution time, and let $N$ be the number of time steps, so that $\Delta t = \frac{T}{N}$ by construction. The first order Trotterization is given by 
$$e^{-i (A + B) \Delta t} \approx e^{-i A \Delta t} e^{-i B \Delta t}$$
where $A$ and $B$ are the two parts of the Hamiltonian. Higher order Trotterization can be achieved by using more complex sequences of exponentials, such as the second order Trotterization:
$$e^{-i (A + B) \Delta t} \approx e^{-i A \Delta t} e^{-i B \Delta t} e^{-i A \Delta t}$$ 
The local truncation error is $O(\Delta t^2)$ for the first order Trotterization and $O(\Delta t^3)$ for the second order Trotterization, and the global error is $\mathcal O(\frac{T^2}{N})$ and $\mathcal O(\frac{T^3}{N^2})$ respectively, assuming that $A$ and $B$ do not have explicit time dependence.\\\\
For a time dependent Hamiltonian, we approximate that the Hamiltonian is roughly constant during each time step, and apply the Trotterization to the Hamiltonian at each time $t$.
\subsection{Implementation Details} 
The implementation of AQC was done using \texttt{qiskit}\cite{javadi-abhariQuantumComputingQiskit2024} 2.3.1. The Trotterized evolution was implemented using the \texttt{PauliEvolutionGate}, \texttt{LieTrotter} and 
\texttt{SuzukiTrotter} classes from \texttt{qiskit}. The algorithm was tested on the \texttt{Aer} simulator, and run on IBM's \texttt{ibm\_kingston} device to analyze performance in a noisy environment. The \texttt{mthree} package was used for readout error mitigation. Calculations were done using \texttt{NumPy}\cite{harrisArrayProgrammingNumPy2020}. The code is available at \href{https://github.com/hkbelagali/Adiabatic-Quantum-Computation}{this} repository.
\section{Results}
\bibliography{CMSE890_Project}
\bibliographystyle{unsrt}
\end{multicols*}
\clearpage
\end{document}